\scanpage{100}

This follows indeed \inserted{in the case where $N'\not\simeq N_\infty$, $N''\not\simeq N_\infty$,} from the second convexity theorem for operators, and from the formulas $\rho_1(uv){-}\rho_n(uv)\leq(\rho_1(u)\rho_1(v)){\cdot-}\rho_n(u)\rho_n(v)$ (number , formula (37)). -- In particular, \deleted{if $N=$ \unsure{the|$N$}} \inserted{for} \deleted{$N$} two polar norms $N$ and $N^\circ$, one finds

\margin{\deleted{\unsure{in the case}}}

\margin{Case $N'\simeq N_\infty$ or $N''\simeq N_\infty$\,: if one admits this $u,v$ \uline{compact} one easily reduces to it in the general case, for in any case $N(u)=\lim_P N(\unsure{P_n u|Pu})$ $P$ orthogonal projections of finite rank.}

\[
\boxed{\lVert uv\rVert_1\leq N(u)N^\circ(v)}\tag*{\deleted{\textcircled{17}}}
\]

One also finds \deleted{directly} \deleted{\unsure{One has also obtained} $\boxed{N(uv)\leq\lVert u\rVert N(v),\quad N(uv)\leq\lVert v\rVert N(u)}$}

The dual of $L_N(E,F)$ is identified with a subspace of the dual of \uline{$L_{N_1}(E,F)=L^{\unsure{1|\infty}}(E,F)$}, hence of $L(F,E)$. \margin{\textcolor{red}{$\vert$}} This subspace contains $L^{N^\circ}(F,E)$, and the map $L^{N^\circ}(F,E)\to(L_N(E,F))'$ is an \uline{isometric} \unsure{monomorphism|isomorphism}.

Let us show that it is \uline{onto}. One may suppose that $N\not\simeq N_1$, let us prove that \deleted{\unsure{if}} \inserted{if a} $\Phi\in(L_N(E,F))'$, \uline{defines} \inserted{$v\in L(F,E)$, \uline{one} has $v\in L^{N^\circ}(F,E)$. Now the norm of} \deleted{a \uline{compact} operator from $F$ to $E$. \unsure{Since} $v^*v$ \unsure{the dual|the dual} uses $|\operatorname{Tr}vu|\leq M N(u)$ for $u\in E'\otimes F$, now $w=v^*v\in L(E,E)$ will have the analogous property, let us prove that $w$ is}

$v=U|v|$, it suffices to prove that $|v|\in L^{N^\circ}(E)$\unsure{$E$|$F$}, now $|v|=\unsure{U'|V}v$, hence has the \unsure{$\subset$|same|same} property analogous to $v$ ($|\operatorname{Tr}(vu)|\leq M N(u)$ for $u\in E'\otimes F$, $M$ constant). One is reduced to $F=E$, $v\geq 0$.

$v$ is compact\,: otherwise there would exist a continuous operator $w$ such that \deleted{\unsure{$v$}} $wv=P$, a Hermitian projection of infinite rank, \deleted{\unsure{one would have the}} one would still have $w\in(L_N(E,E))'$, which is impossible, \deleted{\unsure{one supposes a projection} \inserted{\uline{Hermitian}} \unsure{$p_n$ of rank $n$|of rank $n$}, $p_n\leq P$, \unsure{one would have|would have some}} \inserted{(taking} an orthonormal basis in $P$, one sees that one would have \deleted{$\operatorname{Tr}\unsure{p_n}P=$} $1\in(\ell_N)'=\ell^{N^\circ}$, which is not so. Thus

\margin{?}

one will have $v=\sum\rho_i\bar e_i\otimes e_i$, and from \uline{the \unsure{stupid|said|same} formula} one sees that $(\rho_i)\in(\ell_N)'=\ell^{N^\circ}$, whence the conclusion. -- As a \uline{consequence}, one thus sees that $L^{N^\circ}(E,F)$ is \uline{always} complete, hence every \uline{space} $L^N(E,F)$ is complete hence also $L_N(E,F)$

\margin{101}
